



% \subsection{Planning Benchmarks}
% \label{subsec:bench_planning}

% Planning in LLM-based agents refers to the ability to decompose a problem into intermediate steps before producing a final answer. While widely used in multi-agent frameworks such as CrewAI~\cite{crewai2025}, the evaluation of planning remains fragmented across datasets that target narrow aspects of step-by-step reasoning, numerical computation, or structured decision-making. To provide a unified view of this capability, we construct a planning benchmark suite composed of three widely studied datasets—\emph{GSM8K}~\cite{cobbe2021training}, \emph{CommonsenseQA (CSQA)}~\cite{talmor-etal-2019-commonsenseqa}, and a carefully selected 100-sample subset of \emph{MATH}~\cite{hendrycks2measuring}. These datasets collectively capture short-chain arithmetic reasoning, multiple-choice commonsense inference, and symbolic mathematical problem solving, thereby enabling a consolidated assessment of how LLMs benefit from or are hindered by planning.

% Each dataset in the suite stresses a different facet of planning. \emph{GSM8K} requires short-step numerical reasoning, making it sensitive to whether the plan generation process meaningfully scaffolds the computation. \emph{CSQA} emphasizes symbolic elimination and structured choice among alternatives, enabling analysis of how planning influences multi-step decision-making rather than arithmetic derivation. \emph{MATH-100} introduces more complex algebraic and geometric reasoning, where multi-step symbolic manipulation is required and deviations from expected plan structures can cause extraction failures or incorrect reasoning paths. By combining these datasets under a uniform planning interface, the benchmark provides a coherent lens for analyzing the strengths and limitations of planning mechanisms in LLM-based agents.

% Together, the three datasets and three execution modes constitute a unified planning benchmark that captures a broad spectrum of reasoning behaviors. The benchmark isolates planning as a capability independent of model size, instruction tuning, or prompting style, enabling systematic comparison of structured versus unstructured plan generation and clarifying when planning helps, when it harms, and why certain implementations introduce failure modes unrelated to the underlying model’s reasoning ability.








% \begin{table}[t]
% \centering
% \footnotesize
% \caption{Overview of Datasets Used in Planning Experiment.}
% \vspace{-3mm}
% \label{tab:benchmark_stats_multicol}
% \begin{tabular}{l r | rrrr | rrrr}
% \toprule
% & & \multicolumn{4}{c|}{\textbf{Question Length (words)}} &
%     \multicolumn{4}{c}{\textbf{Gold Answer Length (words)}} \\
% \cmidrule(lr){3-6} \cmidrule(lr){7-10}
% \textbf{Benchmark} & \textbf{\#Q} &
% \textbf{Min} & \textbf{Max} & \textbf{Avg} & \textbf{Total} &
% \textbf{Min} & \textbf{Max} & \textbf{Avg} & \textbf{Total} \\
% \midrule
% GSM8K & 1319 & 15 & 164 & 46.3 & 61,005 & 1 & 1 & 1.0 & 1,319 \\
% CSQA & 1221 & 19 & 69 & 31.6 & 38,615 & 1 & 1 & 1.0 & 1,221 \\
% MATH & 5000 & 2 & 252 & 30.7 & 153,261 & 1 & 19 & 1.3 & 6,345 \\
% \bottomrule
% \end{tabular}
% \end{table}



% The statistics in Table \ref{tab:benchmark_stats_multicol} highlight the structural differences among the three benchmarks used in the planning experiment. GSM8K contains moderately long natural-language questions, averaging 46.3 words, reflecting the multi-step reasoning typical of grade-school math problems. CSQA questions are shorter on average (31.6 words), consistent with their multiple-choice format and highly templated structure. In contrast, the MATH dataset exhibits more variability, ranging from very concise algebraic prompts to lengthy problem-solving descriptions. Although the average MATH question length is only 30.7 words, the maximum length exceeds 250 words, indicating the presence of multi-paragraph, competition-level problems. The gold answers further reinforce this structural diversity: both GSM8K and CSQA provide single-token labels, while MATH contains multi-step symbolic expressions with gold answers reaching up to 19 words.

% One interesting artifact revealed by the analysis is the extremely small minimum question length in MATH (two words), which at first appears counter-intuitive for a mathematical problem-setting benchmark. Upon inspection, this occurs in questions that contain very short natural-language scaffolding followed by a symbolic expression, such as the example “Simplify $(4a^2)^3$.” Although mathematically meaningful, such prompts contain only two English words (“Simplify” and the expression treated as a token) when processed with a whitespace-based tokenizer. This illustrates an important characteristic of the dataset: unlike GSM8K and CSQA, the MATH corpus blends natural language with LaTeX-style symbolic expressions that compress substantial semantic content into compact forms. Consequently, raw word-count statistics for MATH should be interpreted with caution, as mathematical expressions encode complexity not captured by simple tokenization.



% \paragraph{Planning Configuration and Execution Control.}
% \textbf{Planning Cost Profile.} Evaluating planning in LLM-based agents introduces additional computational overhead beyond direct answer generation, as planning-enabled executions require explicit intermediate reasoning steps and, in some cases, separate judging models for structured or numerical correctness. When planning is enabled, agents invoke an auxiliary planning model to generate intermediate decompositions prior to final answer production, increasing both latency and the number of LLM calls per question. In benchmarks such as MATH, this overhead is further amplified by the need for specialized judgment of mathematical correctness, which is not reliably captured by surface-form matching alone. 
% \textbf{Exposed Planning Controls.} To ensure controlled and reproducible evaluation, we expose planning as an explicit configuration toggle that enables or disables intermediate plan generation while keeping all other execution components fixed. The configuration further decouples the backbone inference model, the planning model, and the mathematical judgment model, allowing each role to be instantiated with different LLMs depending on accuracy, cost, or availability constraints. This separation enables systematic analysis of whether gains arise from improved planning quality, stronger base reasoning, or more reliable evaluation, while avoiding confounding effects caused by monolithic model choices. Dataset-level controls additionally allow limiting the number of evaluated instances per benchmark, enabling scalable experimentation from partial diagnostic runs to full benchmark evaluations under identical execution semantics. 




\vspace{-1.5mm}
\subsection{Planning Benchmarks}
\label{subsec:bench_planning}

Planning in LLM-based agents refers to generating intermediate solution steps before producing an answer. While some frameworks, such as CrewAI~\cite{crewai2025}, expose explicit planning, its evaluation remains fragmented across isolated reasoning benchmarks. We therefore adopt \emph{GSM8K}~\cite{cobbe2021training}, \emph{CommonsenseQA (CSQA)}~\cite{talmor-etal-2019-commonsenseqa}, and \emph{MATH}~\cite{hendrycks2measuring} (Table~\ref{tab:benchmark_stats_multicol}), which cover numerical, commonsense, and symbolic reasoning tasks where multi-step planning influences outcomes.\footnote{\url{https://huggingface.co/datasets/openai/gsm8k} (test only); \url{https://huggingface.co/datasets/tau/commonsense_qa} (validation only); \url{https://github.com/hendrycks/math}}

\noindent\textbf{Benchmark Structure and Planning Interfaces.}
Each benchmark is executed under three controlled planning interfaces. \emph{NoPlan} passes tasks directly to the LLM without intermediate reasoning. \emph{Crew-Plan} follows CrewAI’s schema-constrained two-stage interface~\cite{crewai2025}, in which the LLM first produces a structured plan. \emph{Direct-LLM-Plan} allows unconstrained natural-language plan generation injected into the context before answer generation. This setup isolates the effect of planning interface constraints.

\noindent\textbf{Evaluation Setup.}
MAFBench provides a unified planning evaluation pipeline that treats planning as a configurable execution stage while keeping models, prompts, and scoring fixed. The pipeline standardizes dataset loading for GSM8K, CSQA, and MATH, and supports complexity-preserving subsampling for large benchmarks to control question count while maintaining difficulty. In our experiments, we use a 100-problem MATH subset. Formatting failures are tracked separately from reasoning errors. Model selection, planning modes, and run budgets are centrally configured to ensure reproducibility and cost control. All three modes are evaluated across multiple LLM backends, enabling systematic analysis of planning benefits and interface effects.




\begin{table}[t]\centering
\caption{MemoryAgentBench split statistics showing session counts, question distributions, and context lengths (in thousands of words) for each memory competency (AR, TTL, LRU, SF).}
\vspace{-3mm}
\label{tab:memory-bench}
\footnotesize
\begin{tabular}{l r r | rrrr | rrrr}
\toprule
 & & & \multicolumn{4}{c|}{\textbf{Questions}} & \multicolumn{4}{c}{\textbf{Context (k words)}} \\
Split & Sessions & Subtasks & Min & Max & Avg & Total & Min & Max & Avg & Total \\
\midrule
AR & 22 & 4 & 60 & 100 & 90.9 & 2,000 & 50 & 560 & 206 & 4,551 \\
TTL & 6 & 2 & 100 & 200 & 116.7 & 700 & 69 & 1,040 & 236 & 1,417 \\
LRU & 110 & 2 & 1 & 10 & 1.6 & 171 & 48 & 560 & 129 & 14,281 \\
SF & 8 & 2 & 100 & 100 & 100.0 & 800 & 4 & 183 & 63 & 511 \\
\bottomrule
\end{tabular}
\vspace{-2mm}
\end{table}


\begin{table}[t]
\centering
\footnotesize
\caption{Dataset statistics for planning benchmarks, reporting question and gold answer length distributions.}
\vspace{-3mm}
\label{tab:benchmark_stats_multicol}
\begin{tabular}{l r | rrrr | rrrr}
\toprule
& & \multicolumn{4}{c|}{\textbf{Question Length (words)}} &
    \multicolumn{4}{c}{\textbf{Gold Answer Length (words)}} \\
\cmidrule(lr){3-6} \cmidrule(lr){7-10}
\textbf{Benchmark} & \textbf{\#Q} &
\textbf{Min} & \textbf{Max} & \textbf{Avg} & \textbf{Total} &
\textbf{Min} & \textbf{Max} & \textbf{Avg} & \textbf{Total} \\
\midrule
GSM8K~\cite{cobbe2021training} & 1319 & 15 & 164 & 46.3 & 61,005 & 1 & 1 & 1.0 & 1,319 \\
CSQA~\cite{talmor-etal-2019-commonsenseqa} & 1221 & 19 & 69 & 31.6 & 38,615 & 1 & 1 & 1.0 & 1,221 \\
MATH~\cite{hendrycks2measuring}  & 5000 & 2 & 252 & 30.7 & 153,261 & 1 & 19 & 1.3 & 6,345 \\
\bottomrule
\end{tabular}
\vspace{-2mm}
\end{table}
